\documentclass[10pt,a4paper]{article}

% ----------------- Paquetes -------------------%
\usepackage[T1]{fontenc}
\usepackage[utf8]{inputenc}
\usepackage[a4paper,margin=2.5cm]{geometry}
\usepackage{mathptmx}             
\usepackage{changepage} 
\usepackage{setspace}
\usepackage[spanish]{babel}
\usepackage[labelfont=bf,labelsep=space]{caption}
\usepackage{amssymb}
\usepackage{amsmath}
\usepackage{ebgaramond}
\usepackage{placeins}
\usepackage{etoolbox}
\usepackage{setspace}
\usepackage{parskip} 
\usepackage{tcolorbox}
\usepackage{needspace}
\usepackage{xparse}
\usepackage{ocgx2}
\usepackage{hyperref}
\usepackage{hypcap}


%%%%%%%%%%%%%%%%%%%%%%%%%%%%%%%%%%%%%%%%%%%%%%%%%%%%%%%%%%%%%%%%
% --- Fija primero tus valores globales "buenos" ---
\setstretch{1.2}                 % Interlineado global
\setlength{\parskip}{0.7\baselineskip}
\setlength{\parindent}{0pt}
\newlength{\GlobalParskip}
\setlength{\GlobalParskip}{\parskip}

\newcounter{eqnum}[subsection]
\renewcommand{\theeqnum}{\arabic{eqnum}}
\newcounter{alphaitem}
\newcounter{numitem}

%% ------------------- Recuadros----------------------%%
\newcommand{\puntosclave}[1]{%
  \begin{tcolorbox}[
    colframe=black,
    title={Puntos clave},
    before upper={\setlength{\parindent}{0.5cm}\setlength{\parskip}{0.5\baselineskip}}
  ]
  #1
  \end{tcolorbox}
}

\newcommand{\modificaciones}[1]{
  \begin{tcolorbox}[
    colback=white,
    colframe=red,
    title={Modificaciones a realizar},
    before upper={\setlength{\parindent}{0.5cm}\setlength{\parskip}{0.5\baselineskip}}
  ]
  #1
  \end{tcolorbox}
}


%% ------------------- Preguntas TEST----------------------%%
% Opciones: radio (single-choice)
\NewDocumentCommand{\OptRadio}{m m m}{%
  \noindent\CheckBox[radio,name=#1,value=#2,width=1.2ex,height=1.2ex]{}%
  \;\textbf{#2.}~#3\par
}

% Opciones: checkbox (multi-choice)
\NewDocumentCommand{\OptCheck}{m m}{%
  \noindent\CheckBox[name=#1,width=1.2ex,height=1.2ex,bordercolor={0 0 0}]{}%
  \;#2\par
}

% Pregunta single-choice (radio)
% Args: 1=id, 2=enunciado, 3=opciones, 4=solución+justificación, 5=imagen (o {}), 6=origen
\NewDocumentCommand{\PreguntaTestSingle}{m m m m m m}{%
  \par\medskip
  \noindent\textbf{#1.}~#2\par\smallskip

    % Imagen (si no es {})
    \IfBlankTF{#5}{}{%
      \begin{center}
        \includegraphics[width=0.75\linewidth]{#5}
      \end{center}
    }

    \begin{Form}
      #3
    \end{Form}

    \par\smallskip
    \switchocg{sol-#1}{\fbox{Mostrar / ocultar solución}}%
    \hfill{\footnotesize #6}\par\smallskip

    \begin{ocg}{Solución}{sol-#1}{0}
      \noindent #4
    \end{ocg}
  \par\medskip
}

% Pregunta multi-choice (checkboxes)
\NewDocumentCommand{\PreguntaTestMulti}{m m m m m m}{%
  \par\medskip
  \noindent\textbf{#1.}~#2\par\smallskip

    % Imagen (si no es {})
    \IfBlankTF{#5}{}{%
      \begin{center}
        \includegraphics[width=0.75\linewidth]{#5}
      \end{center}
    }
    \begin{Form}
      #3
    \end{Form}

    \par\smallskip
    \switchocg{sol-#1}{\fbox{Mostrar / ocultar solución}}%
    \hfill{\footnotesize #6}\par\smallskip

    \begin{ocg}{Solución}{sol-#1}{0}
      \noindent #4
    \end{ocg}
  \par\medskip
}

%% ------------------- Listas----------------------%%
    % --- Lista alfabética ---
    \newenvironment{la}
      {\begin{list}{\alph{alphaitem})}{%
          \usecounter{alphaitem}%
          \setlength{\leftmargin}{1cm}%
          \setlength{\parskip}{3pt}% 
          \setlength{\itemsep}{0pt}% ← espacio entre ítems
          \setlength{\parsep}{0pt}%  ← párrafos dentro de un ítem
      }}
      {\end{list}}
      
    % --- Lista numérica ---
    \newenvironment{lnum}
      {\begin{list}{\arabic{numitem}.}{%
          \usecounter{numitem}%
          \setlength{\leftmargin}{1cm}%
          \setlength{\parskip}{3pt}% 
          \setlength{\itemsep}{0pt}% ← espacio entre ítems
          \setlength{\parsep}{0pt}%  ← párrafos dentro de un ítem
      }}
      {\end{list}}
      
% ------------------- Imágenes----------------------%
    \graphicspath{{Img/}}% Rutas por defecto 
    % --- Imagen adaptativa ---
    % Registros de dimensión definidos una sola vez
    \newdimen\imgcap
    \newdimen\imgmin
    \newdimen\imgmax
    \newdimen\imgremain
    \newdimen\imgh
    \newdimen\imgneed
    
    \newcommand{\imagenfit}[3]{%
      \addvspace{10pt}% separación antes
      \par\begingroup
        % Parámetros de composición (asignaciones, no \newdimen)
        \imgcap=1\baselineskip            % alto estimado de título+fuente
        \imgmin=0.15\textheight
        \imgmax=0.15\textheight
        \imgremain=\pagegoal
        \ifdim\pagegoal=\maxdimen \imgremain=0pt\fi
        \advance\imgremain by -\pagetotal
        \advance\imgremain by -\imgcap
        % Decisión de altura destino
        \ifdim\imgremain<\imgmin
          \newpage
          \imgh=0.20\textheight
        \else
          \imgh=\imgremain
          \ifdim\imgh>\imgmax \imgh=\imgmax \fi
        \fi
        % Reservar espacio para evitar cortes del bloque completo
        \imgneed=\imgh \advance\imgneed by \imgcap
        \Needspace{\imgneed}
        % Cuerpo
        \noindent\begin{minipage}{\textwidth}
          \centering
          \captionof{figure}{\textemdash\ #2}\par\vspace{0.3em}% título arriba
          \includegraphics[width=\linewidth,height=\imgh,keepaspectratio]{\detokenize{#1}}\par
          \vspace{0.3em}\small\textit{Fuente: }#3% fuente abajo
        \end{minipage}%
      \endgroup
      \par\addvspace{10pt}%
    }

%------------ Bloque de RECUADRO ESPECIAL -------------%
    \newenvironment{specialblock}[1]{%
      \begin{quote} % crea sangría a izquierda y derecha
      \small % texto más pequeño
      \noindent\textbf{#1}\\[-0.3em] % título en negrita
      \rule{\linewidth}{0.4pt}\\[0.6em]}{
      \end{quote}}
      
%-------------------------- Portada -----------------------%
\newcommand{\oldtitle}[1]{\def\OldTitle{#1}}
\newcommand{\changerationale}[1]{\def\ChangeWhy{#1}}

\newcommand{\changebox}{%
\begin{tcolorbox}[title={Historial de cambios del tema}]
\textbf{Título anterior:} \OldTitle\par
\textbf{Motivo del cambio:} \ChangeWhy
\end{tcolorbox}
}

%-------------------------- Author Font -----------------------%
    \newcommand{\authorfont}[1]{{\fontfamily{EBGaramond-TLF}\selectfont #1}}

%-------------------------- Anexos -----------------------%
    \makeatletter
    \newcounter{anexo}
    \renewcommand{\theanexo}{\Roman{anexo}}
    \newcommand{\anexo}[1]{%
      \clearpage
      \phantomsection
      \refstepcounter{anexo}%
      \noindent\textbf{Anexo \theanexo.- }#1\par\medskip
      \addcontentsline{stc}{section}{Anexo \theanexo.- #1}%
    }
    \makeatother

    %-------------------------- Lista de figuras (personal) -----------------------%
\makeatletter
\newcommand{\MyListOfFigures}{%
  \clearpage
  \phantomsection
  {\centering\bfseries\Large Índice de figuras\par}%
  \medskip
  % Entrada en nuestro índice de contenido (stc)
  \addcontentsline{stc}{section}{Índice de figuras}%
  % Recuperación del listado (sin cabecera propia)
  \begingroup
    \parindent=0pt\parskip=2pt
    \@starttoc{lof}%
  \endgroup
}
\makeatother

%-------------------------- Bibliografía -----------------------%
\makeatletter
% Contador para las referencias
\newcounter{myref}
% Cabecera de la bibliografía, entra en el índice 'stc'
\newcommand{\MyBibliography}{%
  \clearpage
  \phantomsection
  {\centering\bfseries\Large Bibliografía y referencias\par}%
  \medskip
  \addcontentsline{stc}{section}{Bibliografía y referencias}%
  \setcounter{myref}{0}%
}
\newcommand{\MyBibItem}[4]{%
  \refstepcounter{myref}%
  \noindent[\themyref]\ #1.\ \emph{#2}. #3. #4\par
}
\makeatother

%--------- Establecer cuatro niveles de índice --------%
    \setcounter{tocdepth}{4}   
    \setcounter{secnumdepth}{4}% numera hasta \paragraph

%%%%%%%%%%%%%%%%%%%%%%%%%%%%%%%%

\makeatletter

    %--------- Interlineado Global --------%
    \edef\GlobalBaseLineStretch{\baselinestretch}
    
    %--------- Espaciado de seccinones --------%
    \renewcommand\section{\@startsection{section}
      {1}{\z@}
      {2.5ex \@plus 1ex \@minus .2ex} % espacio antes
      {1.5ex \@plus .2ex}             % espacio después
      {\normalfont\Large\bfseries}    % formato del título
    }
    \renewcommand\subsection{\@startsection{subsection}{2}{\z@}
      {3ex \@plus 0.8ex \@minus .2ex}
      {1.5ex \@plus .2ex}
      {\normalfont\large\bfseries}
    }
    \renewcommand\subsubsection{\@startsection{subsubsection}{3}{\z@}
      {3ex \@plus 0.5ex \@minus .2ex}
      {1ex \@plus .2ex}
      {\normalfont\normalsize\bfseries}
    }
    \renewcommand\paragraph{\@startsection{paragraph}{4}{\z@}%
    {-2ex \@plus -0.5ex \@minus -.2ex}% espacio antes
    {0.8ex \@plus .2ex}% espacio después
    {\normalfont\normalsize\itshape}}% ← cursiva en lugar de negrita

    %--------- Ecuaciones --------%   
    % Variables globales para almacenar títulos y notas
    \gdef\leq@current@section{}%
    \gdef\leq@current@subsection{}%
    \gdef\leq@last@printed@section{}%
    \gdef\leq@last@printed@subsection{}%
    
    % Comando para añadir nota (invisible en el cuerpo)
    \newcommand{\eqnote}[1]{%
      \addcontentsline{leq}{leqnote}{#1}%
    }
    
    % Comando para ecuaciones
    \newcommand{\eqblock}[2]{%
      % Verificar si necesitamos imprimir el título de sección
      \ifx\leq@current@section\leq@last@printed@section\else
        \addcontentsline{leq}{leqsection}{\leq@current@section}%
        \global\let\leq@last@printed@section\leq@current@section
        \gdef\leq@last@printed@subsection{}% Reset subsección al cambiar sección
      \fi
      % Verificar si necesitamos imprimir el título de subsección
      \ifx\leq@current@subsection\leq@last@printed@subsection\else
        \ifx\leq@current@subsection\@empty\else
          \addcontentsline{leq}{leqsubsection}{\leq@current@subsection}%
          \global\let\leq@last@printed@subsection\leq@current@subsection
        \fi
      \fi
      % Ecuación
      \refstepcounter{eqnum}%
      \begin{equation*}
        #1\tag{E.\theeqnum}
      \end{equation*}%
      \addcontentsline{leq}{leqt}{[E.\theeqnum] -- #2}%
      \addcontentsline{leq}{leqm}{#1}%
    }
    
    %%% ----- Índice de ecuaciones ----------%%%
    % Tamaño para []
    \newcommand{\leqIndexFont}{\normalsize}
    \newcommand{\leqTitleFont}{\tiny}
    \newcommand{\leqNoteFont}{\tiny}
    % Cabecera de sección 
    \newcommand*\l@leqsection[2]{%
      \addvspace{25pt}% Mayor separación antes de nueva sección
      \noindent\leqIndexFont
      \makebox[\linewidth]{\rule{.38\linewidth}{0.4pt}\ \ \textbf{  \MakeUppercase{#1}}\ \ \rule{.38\linewidth}{0.4pt}}%
      \par\vspace{-10pt}
    }
    % Cabecera de subsección 
    \newcommand*\l@leqsubsection[2]{%
      \par\addvspace{15pt}%
      \noindent\leqIndexFont #1
      \par\vspace{-7pt}
    }
    % Título de ecuación 
    \newcommand*\l@leqt[2]{%
    \par\addvspace{10pt}%
      \par\noindent\leqTitleFont #1\par
    }
    % Miniatura de ecuación
    \newcommand*\l@leqm[2]{%
      \vspace{0.3\baselineskip}%
      \par\scriptsize\noindent\hspace*{1.5em}\(#1\)\par%
    }
    % Nota de ecuación 
    \newcommand*\l@leqnote[2]{%
      \vspace{0.2\baselineskip}%
      \par\noindent\leqNoteFont\hspace*{1.5em}\textbf{Nota.--} #1\par\vspace{0.5\baselineskip}%
    }
    % Generación del índice
    \newcommand{\mylistofequations}{%
      \clearpage
      \phantomsection
      % Título visual de la lista (ajústalo a tu gusto):
      {\centering\bfseries\Large Índice de ecuaciones\par}
      % Entrada en nuestro índice 'stc':
      \addcontentsline{stc}{section}{Índice de ecuaciones}%
      % Llamada a tu lista real (usa el nombre exacto de tu macro):
      \@starttoc{leq}%
    }
    
    %--------- Captura de títulos de sección --------%
    \let\leq@original@section\section
    \let\leq@original@subsection\subsection
    % Flag para saber si estamos en una sección asterisco
    \newif\ifLeqStarSection
    \LeqStarSectionfalse
    
    % Redefinir \section CON CONTROL DE ASTERISCO
    \renewcommand{\section}{%
      \@ifstar{\leq@section@star}{\@dblarg\leq@section@nostar}%
    }
    \def\leq@section@star#1{%
      \global\LeqStarSectiontrue
      \gdef\leq@current@section{#1}%
      \gdef\leq@current@subsection{}%
      \leq@original@section*{#1}%
      \global\LeqStarSectionfalse
    }
    \def\leq@section@nostar[#1]#2{%
      \global\LeqStarSectionfalse
      \gdef\leq@current@section{#2}%
      \gdef\leq@current@subsection{}%
      \leq@original@section[#1]{#2}%
    }
    % Redefinir \subsection
    \renewcommand{\subsection}{%
      \@ifstar{\leq@subsection@star}{\@dblarg\leq@subsection@nostar}%
    }
    \def\leq@subsection@star#1{%
      \global\LeqStarSectiontrue
      \gdef\leq@current@subsection{#1}%
      \leq@original@subsection*{#1}%
      \global\LeqStarSectionfalse
    }
    \def\leq@subsection@nostar[#1]#2{%
      \global\LeqStarSectionfalse
      \gdef\leq@current@subsection{#2}%
      \leq@original@subsection[#1]{#2}%
    }

    % ----- Restauración de valores Globales de estilo ----------%
    % Flags
    \newif\ifInFootnote
    \InFootnotefalse
    \pretocmd{\@footnotetext}{\InFootnotetrue}{}{}
    \apptocmd{\@footnotetext}{\InFootnotefalse}{}{}
    % Profundidad de listas (soporta anidadas)
    \newcount\MyListDepth \MyListDepth=0
    \AtBeginEnvironment{la}{\advance\MyListDepth by 1}
    \AfterEndEnvironment{la}{\advance\MyListDepth by -1}
    \AtBeginEnvironment{lnum}{\advance\MyListDepth by 1}
    \AfterEndEnvironment{lnum}{\advance\MyListDepth by -1}
    \AtBeginEnvironment{itemize}{\advance\MyListDepth by 1}
    \AfterEndEnvironment{itemize}{\advance\MyListDepth by -1}
    \AtBeginEnvironment{enumerate}{\advance\MyListDepth by 1}
    \AfterEndEnvironment{enumerate}{\advance\MyListDepth by -1}
    % Restauración diferida: hazlo al INICIO del próximo párrafo real
    \newif\if@pendingRestore
    \@pendingRestorefalse
    \newcommand{\RequestRestore}{\global\@pendingRestoretrue}
    \everypar{%
      \if@pendingRestore
        \global\@pendingRestorefalse
        \ifInFootnote\else
          \ifnum\MyListDepth=0
            \setlength{\parskip}{\GlobalParskip}%
            \setstretch{\GlobalBaseLineStretch}%
          \fi
        \fi
      \fi
    }
    % Al final de bloques:
    \AfterEndEnvironment{equation}{\RequestRestore}
    \AfterEndEnvironment{equation*}{\RequestRestore}
    \AfterEndEnvironment{la}{\RequestRestore}
    \AfterEndEnvironment{lnum}{\RequestRestore}
    % Al final de macros:
    \apptocmd{\eqblock}{\RequestRestore}{}{}
    \apptocmd{\imagenfit}{\RequestRestore}{}{}
    
\makeatother

% ===================== ToC propio con 4 niveles (único bloque) =====================
\makeatletter

% --- Escritores: añadimos una línea al .stc en cada cabecera numerada ---
% Guardamos las versiones actuales para llamar al comportamiento normal
\let\MyTOC@orig@section\section
\let\MyTOC@orig@subsection\subsection
\let\MyTOC@orig@subsubsection\subsubsection
\let\MyTOC@orig@paragraph\paragraph

% SECTION
\renewcommand{\section}{%
  \@ifstar{\MyTOC@section@star}{\@dblarg\MyTOC@section@nostar}%
}
\def\MyTOC@section@star#1{%
  \MyTOC@orig@section*{#1}% sin entrada al .stc
}
\def\MyTOC@section@nostar[#1]#2{%
  \MyTOC@orig@section[{#1}]{#2}% comportamiento normal (escribe al .toc)
  % Escribimos también nuestra línea al .stc (texto con \numberline)
  \addcontentsline{stc}{section}{\protect\numberline{\thesection}#2}%
}

% SUBSECTION
\renewcommand{\subsection}{%
  \@ifstar{\MyTOC@subsection@star}{\@dblarg\MyTOC@subsection@nostar}%
}
\def\MyTOC@subsection@star#1{%
  \MyTOC@orig@subsection*{#1}%
}
\def\MyTOC@subsection@nostar[#1]#2{%
  \MyTOC@orig@subsection[{#1}]{#2}%
  \addcontentsline{stc}{subsection}{\protect\numberline{\thesubsection}#2}%
}

% SUBSUBSECTION
\renewcommand{\subsubsection}{%
  \@ifstar{\MyTOC@subsubsection@star}{\@dblarg\MyTOC@subsubsection@nostar}%
}
\def\MyTOC@subsubsection@star#1{%
  \MyTOC@orig@subsubsection*{#1}%
}
\def\MyTOC@subsubsection@nostar[#1]#2{%
  \MyTOC@orig@subsubsection[{#1}]{#2}%
  \addcontentsline{stc}{subsubsection}{\protect\numberline{\thesubsubsection}#2}%
}

% PARAGRAPH
\renewcommand{\paragraph}{%
  \@ifstar{\MyTOC@paragraph@star}{\@dblarg\MyTOC@paragraph@nostar}%
}
\def\MyTOC@paragraph@star#1{%
  \MyTOC@orig@paragraph*{#1}%
}
\def\MyTOC@paragraph@nostar[#1]#2{%
  \MyTOC@orig@paragraph[{#1}]{#2}%
  % Si no numeras los \paragraph, cambia \theparagraph por texto plano sin \numberline
  \addcontentsline{stc}{paragraph}{\protect\numberline{\theparagraph}#2}%
}

% --- Impresor del índice propio (lee el .stc) ---
\newcommand{\MyTOC}{%
  \par\bigskip
  {\centering\bfseries\Large Índice de contenido\par}%
  \medskip
  \begingroup
    \parindent=0pt\parskip=2pt
    \@starttoc{stc}%
  \endgroup
}

\makeatother
% ===================== fin ToC propio =====================


%%%%%%%%%%%%%%


% --- Datos del tema ---
\title{Tema 3.A.31 -- Análisis de las tablas Input-Output\\
\large }
\author{Marcelo Ortega}
\date{Diciembre 2025}
\newcommand{\TemaCode}{3A31}

\oldtitle{Tema 3.A.27. Análisis de las tablas Input-Output}
\changerationale{Sin cambios.}

\usepackage{soul}\sethlcolor{yellow}% resaltado amarillo: \hl{texto} (Panel TCEE, ⌃H)
\begin{document}


\maketitle
\changebox

\newpage

% --- Página de resumen y modificaciones ---
\puntosclave{ %% Escribir puntos clave Apuntes CECO
    Respecto a la distribución del tiempo:
    \begin{lnum}
        \item El primer bloque (que debería consumir unos 10 minutos de los 30 disponibles para la exposición del tema) consiste en la descripción detallada de la estructura de las tablas input-output.
        \item El segundo bloque consiste en describir una serie de herramientas de simulación que pueden obtenerse cuando se añaden a las tablas input-output un conjunto de hipótesis auxiliares. En este ámbito lo más relevante es tener una buena comprensión de cómo funcionan los efectos multiplicadores que se derivan de la existencia de coeficientes técnicos no nulos y de funciones de producción de coeficientes fijos.
    \end{lnum}
    
    }
\vspace{1cm}

\modificaciones{
    
    }

\newpage
\MyTOC
\newpage

\mylistofequations
\clearpage

\MyListOfFigures
\clearpage


\pagenumbering{arabic}
\setcounter{page}{1}


\section{Introducción}\label{intro}


\subsection{Contextualización}\label{context}


\subsubsection{Relevancia}\label{importance}
Es relevante el análisis input-output porque proporciona una visión de una economía más pormenorizada que cualquier otro instrumento y además porque pem1ite llevar a cabo simulaciones de incidencia en el resto del sistema económico -en términos de producción, exportaciones, precios y empleo- de cualquier variación en una actividad económica. Finalmente el sistema input-output refleja muy bien la idea de interdependencia, una idea central a la economía.

\subsubsection{Historia}\label{history} 
Toda ciencia se inicia con estudios parciales y locales hasta que aparece un problema fundamental o una gran sintesis. Filósofos, arbitristas, moralistas y otros estudiosos, así como hombres de negocios, habían estudiado cuestiones concretas como los impuestos, el comercio o el dinero, pero QUESNAY, un médico familiarizado con la circulación de la sangre, avanzó la idea del flujo circular del producto social y a partir de ahí surge la idea de la interdependencia de los fenómenos económicos.

WALRAS lleva la idea de la interdependencia de los fenómenos económicos a su desarrollo más acabado, elaborando un modelo matemático de equilibrio general, mediante un sistema de ecuaciones simultáneas, con el que describía la determinación simultánea de los precios y las cantidades en un sistema económico.\footnote{%
    En su monumental Historia del \textit{Análisis Económico} J.A. SCHUMPETER considera la obra de WALRAS el trabajo teórico de más categoría realizado por un economista (Pág. 1113), denominando a su obra la piedra miliar y la carta magna de la economía.
}

Leontief conecta con esta gran corriente de interdependencia en Economía, aunque sei'lala que Quesnay se basaba en datos imaginarios. Pero en el siglo XX ya existía un gran volumen de datos primarios, faltaban las cajas de integración de tales datos, una base teórica que orientase su captación y organización con sentido para la ciencia. La pretensión de Leontief es aportar un esquema teórico con capacidad para interpretar la multitud de datos primarios en un esquema coherente con capacidad de explicar y medir las interrelaciones que existen en un sistema económico. En un artículo de divulgación científica para la revista Scientiflc American (1951) escribía que tenemos una alta concentración de teoría sin hechos y, por otra parte, una alta acumulación de hechos sin teoría. Por eso se propone rellenar las empty boxes de la teoría con hechos relevantes.

\subsubsection{Problemática}\label{problem} 


\subsection{Estructura de la exposición}\label{structure}



\clearpage
\section{Las TIO en el SCN}
Las transacciones que se realizan en un sistema económico se describen con eficacia como un circuito económico o sistema de flujos interdependientes, tanto de bienes y servicios como de factores productivos. La tabla input-output pone énfasis en estos flujos desde el punto de vista de las ramas de actividad. 

El amplio desglose de éstas que realiza la TIO en comparación con otros instrumentos contables es su más genuina contribución como documento del sistema de contabilidad nacional, donde se lleva a cabo una desagregación de la economía por ramas de actividad y productos señalando las operaciones de bienes y servicios entre las ramas de actividad y con los demandantes finales. La elaboración de la tabla permite examinar la lógica y coherencia de las cuentas nacionales en un cuadro detallado al incorporar los tres enfoques utilizados para la estimación del PIB (valor añadido, rentas y gasto). 

El marco input-output se compone de tres tablas: la tabla de origen, la tabla de destino y la tabla simétrica. Para desarrollar toda nuestra exposición, supondremos una economía abierta con dos sectores productivos: $S_1$ y $S_2$.

\subsection{La Tabla de Origen}
La tabla de origen muestra la oferta de bienes y servicios de una economía desglosando éstos por productos y ramas de actividad, distinguiendo asimismo entre producción nacional e importaciones. Puede presentarse a precios básicos y precios de adquisición.

\imagenfit{TablaOrigen.png}{Tabla de Origen para dos sectores de la economía y $M$ productos}{Elaboración propia}

Es decir, en términos contables (atendiendo a la SEC-2010),a partir de ella pueden obtenerse los \textbf{recursos} de una economía por productos encuadrados en cada rama de actividad. 

La producción total por ramas más las importaciones dan la oferta total de bienes y servicios a precios básicos. Si se añaden a cada producto los márgenes de comercio y de transporte así como los impuestos netos sobre los productos se obtiene la oferta de cada producto a precios de adquisición.

\subsection{La Tabla de Destino}
La tabla de destino es también una matriz en la que por columnas se presentan las ramas de actividad y por filas los productos. Esta tabla muestra los \textbf{empleos} de los bienes y servicios por productos y tipo de empleo (intermedio y cada uno de los empleos finales).

\imagenfit{TablaDestino.png}{Tabla de Destino para economía de dos sectores y $M$ productos}{Elaboración propia}

Como vemos en la tabla de origen se registran los valores precios de adquisición mientras que en la tabla de destino se registran a precios de adquisición. El V.A. se registra a precios básicos y puede obtenerse como resultado de la producción a precios básicos menos los consumos intermedios a precios de adquisición. Además, en las dos tablas anteriores se producen dos identidades, siempre que la oferta y los empleos se valoren correctamente:
\begin{lnum}
    \item \textit{Por ramas de actividad}. La producción por rama (tabla de origen) es igual a los insumos por rama (tabla de destino), esto es, la producción por rama es igual a los consumos intermedios más el valor añadido; y,
    \item \textit{Por productos}. La oferta total por producto (tabla de origen) es igual a los empleos totales por producto (tabla de destino). 
\end{lnum}

Estas identidades revelan la utilidad de estas tablas desde el punto de vista contable, al permitir contrastar la calidad de las estadísticas por fuentes y puntos de vista diferentes, pudiendo llenar lagunas y cotejar la corrección de los datos. En la tabla de origen los productos se registran a precios básicos. 

\subsection{TIO simétrica}

\subsubsection{La construcción por bloques}
En una TIO pueden distinguirse tres grandes bloques: bloque de relaciones interindustriales o de consumos intermedios, bloque de demandas finales y bloque de inputs primarios:

\imagenfit{EstructuraTIOSimetrica.png}{Estructura de la TIO simétrica}{Apuntes ICEX-CECO. Tema 3.A.30}

\paragraph{Relaciones Interindustriales: Bloque I}
La parte más sustantiva de una TIO, en comparación con otros elementos de la contabilidad nacional, es el bloque de consumos intermedios, que muestra las relaciones interindustriales de cada rama con el resto. A cada rama le corresponde una columna y una fila de la TIO, leyéndose en columnas los inputs o entradas para el proceso productivo de cada rama de productos de otras ramas y, en filas, los outputs , salidas o destinos de las producciones de cada rama. Por tanto, las columnas representan los productos que cada rama utiliza de las demás para obtener su producción, mientras que por filas indica los outupts o destinos de los productos de cada rama que se utilizan como consumos intennedios por parte de otras ramas.

\imagenfit{BloqueI.png}{Relaciones Interindustriales: Bloque I}{Elaboración propia}

Por definición, el total de consumos intermedios utilizados por todas las ramas (suma de los totales por columnas, $j$) coincide con el total de salidas (suma por filas, $i$) de productos para uso intermedio del resto de ramas. Es decir, el total de consumos intermedios coincide por filas y por columnas. Pero esta igualdad no tiene por qué producirse rama a rama; no tiene por qué coincidir el montante de productos que utiliza la agricultura, como semillas, fe rtilizantes, energía y otros, con los productos agrícolas que se destinan a consumos intermedios de otras ramas, como a la industria alimentaria o los restaurantes.


\paragraph{La Demanda Final: Bloque II}
Parte de la producción de las ramas se destina a usos finales (final se opone a intermedio): consumo individual (C), consumo colectivo (CC), formación bruta de capital (FBK) y exportaciones (E). Tales destinos se recogen en el Bloque II de la TIO.

\imagenfit{BloqueII.png}{Demada Final: Bloque II}{Elaboración propia}

Huelga decir que los bienes de capital, aunque se destinen, como es habitual, a trabajos de producción por las empresas adquirientes, no son bienes que se incorporan enteramente a las producciones fabricadas cada año, como ocurre con las materias primas y otros productos intermedios, sino que permanecen durante varios periodos como bienes de capital. Por ello, este destino es un destino final y no intermedio.

\paragraph{Los Inputs Primarios: Bloque III}
Los factores de producc ión primarios, trabajo y capital, se remuneran, respectivamente, con salarios y excedente bruto de explotación. Todo ello se registra en el Bloque III de la TIO.

\imagenfit{BloqueIII.png}{Inputs Primarios (remuneración): Bloque III}{Elaboración propia}

Sumando el VAB de cada rama con el total de sus consumos intennedios, se obtiene el valor de la producción de dicha rama.

\subsubsection{La Tabla I-O}
La agregación de los tres bloques permite tener una tabla input-output simplificada. Distingueremos según se trate de un desglose con sector exterior o no.

\paragraph{TIO SImplificada}
En una TIO agrupada, donde no se representa explícitamente el sector exterior de la economía, por lo que respecta a las importaciones, tendremos una disposición como laque sigue:

\imagenfit{TablaSimplificadaI.png}{Tabla Simplificada (Economía cerrada, no importaciones)}{Elaboración propia}

Puede decirse que una columna de una TIO expresa la estructura productiva de una rama (estructura de costes de un sector productivo), esto es, los consumos intermedios y los inputs primarios utilizados para obtener su producción; mientras que las filas de una TIO representan los destinos (intermedios y finales) a los que van los productos de cada rama, por tanto representaría más bien la estructura del ingreso o de la demanda. 

\paragraph{TIO con importaciones}
La TIO simétrica, que se ha presentado hasta ahora por simplificación didáctica, es la correspondiente a un país que no efectúa importaciones, aunque sí exportaciones. 

Como lo habitual es que se efectúen también importaciones, tanto para ser utilizadas en los procesos productivos (como consumos intermedios), como para los utilizadores finales (bienes de consumo o de capital), la TIO recoge esta circunstancia, añadiendo, en la parte inferior de la columna correspondiente a cada rama, una fila de importaciones (M), esto es, importaciones similares a los productos de cada rama, es decir productos recogidos en los dígitos de la CNAE que han sido incluidos en la rama. Es decir, por la homogeneidad de las filas y las columnas de una TIO, las importaciones deben añadirse a \textbf{la rama que efectúa tales producciones} (producciones similares a los productos objeto de importación), aunque en la realidad las importaciones hayan sido realizadas por empresarios pertenecientes a otra rama.\footnote{%
    Por ejemplo, si los comerciantes efectúan importaciones de ordenadores para vender en el país, tales importaciones se añadirán a la rama de Fabricación de ordenadores y equipo informático y no a la rama de comercio. De otro modo, se perdería la homogeneidad de las ramas.
}

\imagenfit{TIO con Economía abierta.png}{Tabla Input-Output en Economía abierta (con importaciones)}{Elaboración propia}

Las importaciones, a su vez, serán objeto de destinos intennedios o finales ( entre estos últimos puede también darse el caso de importaciones que podrían ser objeto de exportación). Las tablas permiten registrar el origen (interior o importado) de cada consumo intennedio o cada demanda final, suministrando la infonnación desagregada en la casilla correspondiente. Así, cada casilla de una TIO recoge tres cantidades correspondientes al origen del producto: total, interior e importado.

\subsubsection{Representación matemática: Operadores}
Para operar, se requiere una presentación más fomializada del modelo input-output. Ello puede realizarse mediante símbolos:
\begin{lnum}
    \item $x_{ij}$. Es una casilla (cruce de una fila y una columna) del bloque de consumos intermedios de la TIO (utilización que la rama $j$ hace de productos de la rama $i$);
    \item $X_J$. Es la producción de la rama $j$ (suma de consumos intermedios de la rama más el valor añadido); 
    \item $VAB_j$ es el valor añadido por la rama $j$, integrado por remuneración de asalariados ($RA_j$) y excedente bruto de explotación ($EBE_j$);
    \item $D_i$. Representa la demanda final de $i$, y se compone por: consumo individual ($C_i$), consumo colectivo ($CC_i$), formación bruta de capital ($FBK_i$) y exportaciones ($E_i$);
    \item $M_j$. Representa las importaciones; y,
    \item $TO_j$ es el total de recursos (producción más importaciones) y $TE_i$ el total de empleos (demanda intermedia más demanda final).
\end{lnum}

En cada rama se produce, en ausencia de importaciones, un equilibrio entre la oferta y la demanda, entre la producción y sus destinos, tal como se expresa en cada fila.

\eqblock{
\begin{aligned}
&\begin{cases}
    TE_i\sim X_i=\sum_jx_{ij}+D_i=\underbrace{W_i}_{\text{DI de $i$}}+D_i\\[0.5em]
    TR_j\sim X_j=\sum_ix_{ij}+VAB_j=\underbrace{U_j}_{\text{CI de $j$}}+VAB_j
\end{cases}\;\Longrightarrow\;
\begin{cases}
    \sum_iX_i=\sum_i\sum_jx_{ij}+\sum_iD_i\\
    \sum_jX_j=\sum_j\sum_ix_{ij}+\sum_jVAB_j
\end{cases}\\[1em]
&\text{Como, $\sum_iX_i=\sum_jX_j$, entonces:}\qquad
\underbrace{\boxed{\overbrace{\sum_iD_i}^{PIB_{\text{\scriptsize Demanda}}}=\underbrace{\sum_jVAB_j}_{PIB_{\text{\scriptsize Oferta}}}}}_{\text{Sin importaciones}}\\[1em]
&\text{Con importaciones:}\quad \boxed{\sum_iD_i-\sum_iM_i=\sum_jVAB_j}
\end{aligned}
}{TIO. Agregados Macroeconómicos e identidades.}

\paragraph{Los coeficientes técnicos}
Puede representarse una estructura de relaciones entre las ramas de una economía calculando, para cada rama, la proporción de cada input en su producción. A tales proporciones se les denomina \textit{coeficientes de input}, que pueden representarse en una tabla o matriz en la que cada elemento es el correspondiente a la tabla original. 

\begin{specialblock}{Coeficientes técnicos}
    Se representan por a $a_{ij}$ y se definen como la utilización que la rama $j$ hace de productos de la rama $i$ por unidad de producción, es decir: $a_{ij}= x_{ij}/X_j$. 
    
    Mucho ojo porque se definen en términos monetarios (valor de los productos de la rama $j$ respecto al valor total de la producción de la rama $i$), por tanto, no responden a ninguna estructura tecnológica subyacente, simplemente reflejan las transacciones económicas entre sectores sin poder realizar ninguna conclusión en términos de tecnología (esto causará problemas, como veremos más adelante).
\end{specialblock}

Los coeficientes técnicos son un tipo privilegiado de coeficientes de input y expresan la utilización que cualquier rama hace de productos de otra por unidad de producción. Por ejemplo, la utilización que la $S_2$ hace de productos de $S_1$ por unidad de producción será: $a_{12}=x_{12}/X_2$ , esto es, utilización que $S_2$ hace de productos de $S_1$, tal como viene cuantificada en la casilla correspondiente (fila 1, columna 2) de la TIO, dividida por la producción total de $S_2$ ($X_2$). Por tanto, siguiendo con el ejemplo de las TIO propuestas:

\eqblock{
\begin{aligned}
    \left\{\begin{matrix}
        a_{11}&a_{12}\\
        a_{21}&a_{22}
    \end{matrix}\right\}=\left\{\begin{matrix}
        x_{11}&x_{12}\\
        x_{21}&x_{22}
    \end{matrix}\right\}\cdot 
    \left\{\begin{matrix}
        1/X_1&0\\
        0&1/X_2
    \end{matrix}\right\}\to\\[1em]
    \to\left\{\begin{matrix}
        2&1\\
        3&1
    \end{matrix}\right\}\cdot\left\{\begin{matrix}
        1/6&0\\
        0&1/7
    \end{matrix}\right\}=\left\{\begin{matrix}
        0.3&0.14\\
        0.5&0.14
    \end{matrix}\right\}=A
\end{aligned}
}{Derivación de los Coeiciente técnicos de la TIO}

\subsection{Aplicaciones prácticas}


\subsubsection{Eslabonamiento interindustrial}
El concepto de eslabonamientos o encadenamientos entre ramas. Existen actividades como la construcción con fuertes eslabonamientos hacia atrás, es decir, que mediante su demanda de inputs intermedios (cemento, vidrio, metal, ... ) serían capaces de promover la aparición o expansión de las actividades suministradoras de los mismos. Pero también se pueden producir eslabonamientos hacia delante, es decir, vinculaciones que se producen cuando se desarrolla una actividad cuyos productos serán utilizados por otras ramas como inputs intennedios para su proceso de producción. Manteniendo el ejemplo de la construcción, ésta generará por ejemplo este tipo de eslabonamientos sobre los servicios in mobiliarios, ya que una expansión de la construcción contribuye al desarrollo de este tipo de servicios.\footnote{%
    La estrategia de desarrollo des equilibrado, propuesta por Albert O. Hirschman, en su importante libro la estrategia del desarrollo, se f undamenta, precisamente, en la capacidad de las llamadas industrias clave ( o actividades con f uertes eslabonamientos hacia atrás y hacia delante) para producir desequilibrios (severa escasez de inputs o exceso de oferta de los mismos) y provocar, de este modo, la aparición de actividades complementarias, ya que tales desequil ibrios actuarían como señales de mercado revelando oportunidades de inversión y estimulando así la actividad económica.
}

Para cada rama de actividad, dichos eslabonamientos, tanto hacia atrás como hacia delante, se pueden tratar de cuantificar considerando sus relaciones con el conjunto del sistema o con cada una de las demás ramas por separado. En el primer caso, se definen como eslabonamientos globales y se presentan las propuestas de Chenery y Watanabe, y Rasmussen; en el segundo caso hablamos de eslabonamientos específicos y se presenta la propuesta de coeficientes de Streit.

\paragraph{Coeficientes de CHENERY y WATANABE}
CHENERY y WATANABE elaboraron métodos para valorar los eslabonamientos, seleccionando aquellas actividades cuyos efectos de eslabonamiento eran superiores a la media, por su capacidad de incidencia en el sistema económico. Para ello, se ordenan las ramas, por su comparación a la media, mediante la combinación de los dos criterios siguientes:
\begin{la}
    \item Utilización de cada rama de inputs intermedios con respecto a su producción; y,
    \item Destino intermedio de los productos de cada rama con relación al total de los destinos de sus producciones.
\end{la}

Así, para medir los eslabonamientos tanto hacia adelante como hacia atrás de la rama productiva $j$, utilizamos:
\eqblock{
\begin{aligned}
    &\text{Eslabonamiento hacia atrás}:\\[1em]
    &\underbrace{\mu_j=\frac{\sum_ix_{ij}}{X_j}=\frac{\text{Consumos Intermedios}_j}{X_j}}_{\text{Proporción de los \textit{inputs} intermedios sobre $X_j$}}\quad \Longrightarrow\quad 
    \boxed{\bar \mu_j=\frac{\sum_j\mu_j}{\#J}}\to\text{Media observada}\\[2em]
    &\text{Eslabonamiento hacia adelante}:\\[1em]
    &\underbrace{\omega_j=\frac{\sum_ix_{ij}}{X_j}=\frac{\text{Demanda Intermedia}_j}{X_j}}_{\text{Proporción de los \textit{inputs} intermedios sobre $X_j$}}\quad \Longrightarrow\quad 
    \boxed{\bar \omega_j=\frac{\sum_j\omega_j}{\#J}}\to\text{Media observada}\\[2em]
\end{aligned}
}{Coeficientes CHENERY y WATANABE. Eslabonamientos}

A partir de los resultados, la combinación de ambos coeficientes permite clasificar las diferentes ramas de una tabla input-output de acuerdo a las casillas de un cuadro como el que se presenta a continuación:

\imagenfit{ClasificacionCHENERY_Eslabonamientos.png}{Clasificación de las ramas de actividad atendiendo al grado de sus eslabonamientos. CHENERY y WATANABE}{Apuntes ICEX-CECO. Tema 3.A.31}

Como observaciones sobre la utilidad de estos coeficientes pueden señalarse las siguientes: 
\begin{lnum}
    \item Utilizan coeficientes directos, no totales. Estos últimos podrían suministrar información más completa;
    \item Son coeficientes medios sin ponderaciones. Por ello no deben extraerse conclusiones precipitadas, ya que una rama puede tener coeficientes $\mu_j$ y $\omega_j$, elevados pero ser irrelevante desde el punto de vista del empleo, de la demanda final o de cualquier otro criterio de importancia económica; y,
    \item Los autores citados no establecen ningún criterio de desviación de los coeficientes. Es decir, los coeficientes anteriores son medias que no tienen en cuenta si las relaciones de una rama con el resto están muy concentradas, de modo que dos ramas con idéntico coeficiente $μ_j$ pueden tener incidencia económica diferente en función de la concentración, ya que un coeficiente de igual cuantía puede provenir de utilizar productos de una sola rama o de distribuir las compras entre muchas.
\end{lnum}

Estas limitaciones evidentes fueron resueltas, en cierta medida, por la formulación de los coeficientes de RASMUSEN\footnote{
Véase [\authorfont{Ver Apuntes ICEX-CECO Tema 3.A.31}]
}

\paragraph{Coeficientes de STREIT}
Los coeficientes presentados hasta aquí, relacionan cada rama con todas las demás de forma conjunta. 

Sin embargo, las relaciones rama a rama también pueden ser objeto de cuantificación para conocer cuáles aparecen habitualmente como inten-elacionadas, pudiendo llegar a formar complejos industriales y establecerse juntas en el espacio. Para este estudio de relaciones rama a rama puede valorarse qué representa cada rama para otra cualquiera. como suministradora de inputs o como receptora de sus productos. De este modo, pueden establecerse cuatro relaciones o ligazones entre dos ramas.

\textbf{Continuar con APUNTES ICEX-CECO. Tema 3.A.31}

Como acabamos de ver, a través del estudio de las TIO podríamos determinar cuestiones tales como qué ramas de actividad son las que más peso tienen en una economía, cuáles son más intensivas en consumos intermedios, cuáles generan mayor valor añadido, etc. además, este tipo de análisis se puede poner en perspectiva mediante comparaciones internacionales con los datos procedentes de las tablas elaboradas para otros países. 

Ahora bien, si a este tipo de cálculos les añadimos una serie de hipótesis, podremos obtener una herramienta de simulación muy relevante para el análisis económico. Esto es lo que pasamos a examinar a continuación. Concretamente, destacaremos dos tipos de simulaciones posibles:
\begin{la}
    \item Estimación de los niveles de producción de cada rama necesarios para satisfacer un obj etivo de demanda final determinado de forma exógena;
    \item Alteración de los precios de cada rama como resultado de la variación de algún elemento del valor añadido o de los precios de los productos de alguna rama.
\end{la}

Empecemos por el primero.

\clearpage
\section{El modelo de demanda}

A la hora de analizar el modelo de demanda, conviene comenzar enumerando los principales supuestos sobre los que se basa éste:
\begin{lnum}
    \item \textit{Homogeneidad}. Cada mercancía es suministrad a únicamente por una rama;
    \item \textit{Proporcionalidad constante}. La función de producción es homogénea de grado uno (lo que implica que no existen economías de escala);
    \item \textit{Aditividad}. El efecto total de llevar a cabo varios tipos de producción es igual a la suma de los efectos separados (esta hipótesis excluye la existencia de economías o deseconom ías externas);
    \item \textit{No restricción de recursos}. Leontief asume que no hay factores limitados; y,
    \item \textit{Coeficientes fijos}. Los coeficientes de producción son fijos, por lo que no hay posibilidad de sustitución entre inputs intermedios, ni entre estos y los inputs primarios. Es decir, la función de producción es de coeficientes fijos. Esto elimina todos los efectos-precio en la sustitución de inputs, y suprime muchos de los ajustes que se dan en el concepto walrasiano de equilibrio general.
\end{lnum} 

Partiendo de estas hipótesis, procedemos ahora a examinar el modelo y sus principales aplicaciones.






\clearpage
\section{El modelo de precios}

\clearpage
\section{Conclusión} 




\subsection{Recapitulación}



\subsection{Extensiones y relación con otras partes del Temario}



\subsection{Opinión}

\subsubsection{Idea final (Salida o cierra)}

\clearpage
\pagenumbering{gobble}
\MyBibliography
\MyBibItem{Autor}{Título}{Año}{DOI -- LINK}


% --- Anexos (no aparecen en el índice) ---
\clearpage
\anexo{\textbf{Preguntas Test}}
%\input{\TemaCode/questions.tex} 



\clearpage
\anexo{Las Tablas Input-Output en el Sistema de Cuentas Nacionales}
Las transacciones que se realizan en un sistema económico se describen con eficacia como un circuito
económico o sistema de flujos interdependientes, tanto de bienes y servicios como de factores
productivos. La tabla input-output pone énfasis en estos flujos desde el punto de vista de las ramas
de actividad. El amplio desglose de éstas que realiza la TIO en comparación con otros instrumentos
contables es su más genuina contribución como documento del sistema de contabilidad nacional. Las
tablas input-output pueden considerarse como una ampliación de las cuentas nacionales, poniendo
énfasis en las transacciones que tienen lugar entre las diferentes ramas de actividad (química,
alimentaria, etc.) en las que pueden agruparse las actividades que tienen lugar en una economía.
El Sistema europeo de cuentas nacionales y regionales de la Unión Europea elabora el marco
contable, compatible entre los países de la U .E. para describir de una manera sistemática y detallada
las características económicas de un país o región. Para ello se definen los conceptos de una forma
armonizada, coherente y operativa.
En el SEC se definen:
A) Las unidades estadísticas y su agregación
. B) Los flujos
C) El sistema de cuentas
D) El marco input-output.
En el marco input-output se lleva a cabo una desagregación de la economía por ramas de actividad y
productos señalando las operaciones de bienes y servicios entre las ramas de actividad y con los
demandantes finales.
La elaboración de la tabla permite examinar la lógica y coherencia de las cuentas nacionales en un
cuadro detallado al incorporar los tres enfoques utilizados para la estimación del PIB (valor añadido,
rentas y gasto).
El marco input-output se compone de tres tablas: la tabla de origen, la tabla de destino y la tabla
simétrica. Comenzaremos describiendo brevemente las dos primeras, para pasar después a centrarnos
en la tabla si.Jn��trica, por ser ésta la más relevante para fines analíticos.


\end{document}